\subsection{局部凸 Lusin 空间}

命题 1. --- 设 E 是一个 Hausdorff 局部凸空间。假设存在一个由满足第一可数公理的 Fréchet 空间构成的序列 \((\mathrm{E}_{n})_{n\in\mathbb{N}}\) ，以及连续线性映射 \(u_{n}:E_{n}\to E\) ，使得 \(\mathrm{E}=\bigcup_{n\in\mathbb{N}}u_{n}(\mathrm{E}_{n})\) 。则 E 是一个 Lusin 空间。

设 \(P_n\) 是 \(u_n\) 的核；则 \(u_n\) 定义从商空间 \(E_n / P_n\) 到 \(u_n(E_n)\) 上的一个连续双射。由于 \(E_n / P_n\) 是满足第一可数公理的 Fréchet 空间（GT, IX, § 3, No.~1），从而是一个波兰空间（GT, IX, § 6, No.~1, 定义 1），故 \(u_n(E_n)\) 是 E 的一个 Lusin 子空间（GT, IX, § 6, No.~4, 命题 11）。于是由 GT, IX, § 6, No.~7, 定理 3 的推论，作为正则空间（GT, III, § 3, No.~1）的 E 是一个 Lusin 空间。

例 1. --- 每个满足第一可数公理的 Fréchet 空间都是一个波兰空间，从而是一个 Lusin 空间。因此，空间 \(\mathcal{C}(\mathrm{X})\) （其中 \(\mathrm{X}\) 局部紧且具有可数基）亦然（\(\mathcal{C}(\mathrm{X})\) 的拓扑取为紧收敛拓扑，cf.~GT, X, § 3, No.~3, 推论 与 § 1, No.~6, 推论 3）；* 空间 \(\mathcal{C}^{\infty}(\mathrm{U})\) （其中 U 是 \(\mathbf{R}^{n}\) 的开子集（III, 第 9 页））与空间 \(\mathcal{H}(\mathrm{U})\) （其中 U 是 \(\mathbf{C}^{n}\) 的开子集（III, 第 10 页））也都如此。

命题 1 表明，空间 \(\mathcal{C}_0^\infty (\mathbf{U})\) （其中 \(\mathbf{U}\) 是 \(\mathbf{R}^n\) 中的开集）、\(\mathcal{G}_s(\mathbf{I})\) （其中 \(\mathbf{I}\) 是 \(\mathbf{R}\) 中的紧区间且 \(s\geqslant 1\) ）以及 \(\mathcal{H}(\mathbf{K})\) （其中 \(\mathbf{K}\) 是 \(\mathbf{C}^n\) 的紧子集）都是 Lusin 空间（III, 第 10 页）。

定理 2. --- 设 E 是一个局部凸空间，它是 E 的子空间所成的递增序列 \((\mathrm{E}_n)_{n\in \mathbb{N}}\) 的归纳极限，诸子空间赋予满足第一可数公理的 Fréchet 空间拓扑。假设 E 的每个紧子集都含于某个 \(E_{n}\) 之中，且在这个空间中是紧的。设 F 是满足第一可数公理的 Fréchet 空间。则空间 \(\mathcal{L}_{c}(E;F)\) 是一个 Lusin 空间。

空间 E 是有界型空间（III, 第 12 页），因此空间 \(\mathcal{L}_c(\mathrm{E};\mathrm{F})\) 是完备的（III, 第 23 页, 命题 12）。线性映射 \(j\colon f\mapsto (f|\mathrm{E}_n)_{n\in \mathbf{N}}\) 是从 \(\mathcal{L}_c(\mathrm{E};\mathrm{F})\) 到乘积空间 \(\prod_{n\in \mathbf{N}}\mathcal{L}_c(\mathrm{E}_n;\mathrm{F})\) 中的一个单射；凭借关于 E 的紧子集的假设以及诸 \(\mathfrak{S}\) -拓扑的定义，\(j\) 是从 \(\mathcal{L}_c(\mathrm{E};\mathrm{F})\) 到它的像上（该像赋予由乘积拓扑诱导的拓扑）的一个同构；此外，由于 \(\mathcal{L}_c(\mathrm{E};\mathrm{F})\) 完备，这个像是 \(\prod_{n\in \mathbf{N}}\mathcal{L}_c(\mathrm{E}_n;\mathrm{F})\) 的一个闭子空间

（GT, II, § 3, No.~4, 命题 8）。于是由 GT, IX, § 6, No.~4，只需证明诸空间 \(\mathcal{L}_c(\mathrm{E};\mathrm{F})\) 中的每一个都是 Lusin 空间。在证明的其余部分，我们将假设 E 是满足第一可数公理的 Fréchet 空间。

由于 F 是满足第一可数公理的 Fréchet 空间，它同构于 Banach 空间的一个可数乘积的闭子空间，其中每个 Banach 空间 \(F_{n}\) 都是 F 的商（II, 第 5 页），因而满足第一可数公理。线性映射 \(j': f \mapsto (\mathrm{pr}_{n} \circ f)_{n \in \mathbb{N}}\) 是从 \(\mathcal{L}_{c}(\mathrm{E}; \mathrm{F})\) 到乘积空间 \(\prod_{n \in \mathbb{N}} \mathcal{L}_{c}(\mathrm{E}; \mathrm{F}_{n})\) 中的一个单射，并且利用诸 S-拓扑以及乘积中开集的定义，\(j'\) 是从 \(\mathcal{L}_{c}(\mathrm{E}; \mathrm{F})\) 到它的像上的一个同构；此外，由于 \(\mathcal{L}_{c}(\mathrm{E}; \mathrm{F})\) 完备，这个像是 \(\prod_{n \in \mathbb{N}} \mathcal{L}_{c}(\mathrm{E}; \mathrm{F})\) 的一个闭子空间。因此只需证明诸空间 \(\mathcal{L}_{c}(\mathrm{E}; \mathrm{F}_{n})\) 中的每一个都是 Lusin 空间（GT, IX, § 6, No.~4），从而我们可以假设 F 是满足第一可数公理的 Banach 空间。

空间 \(\mathcal{L}_{c}(\mathrm{E};\mathrm{F})\) 是等度连续闭子集所成可数族的并（III, 第 19 页, 推论 1 与 GT, X, § 2, No.~3, 命题 6）。但 \(\mathcal{L}_{c}(\mathrm{E};\mathrm{F})\) 的每个等度连续子集 H 都是可度量化的且满足第一可数公理（III, 第 18 页, 命题 6 与 GT, X, § 2, No.~4, 定理 1）；若 H 是闭的，则它对由 \(\mathcal{L}_{c}(\mathrm{E};\mathrm{F})\) 的一致结构诱导的一致结构而言是完备的空间，因为后者是完备的。换言之，H 是一个波兰空间，从而更是 Lusin 空间；于是正则空间 \(\mathcal{L}_{c}(\mathrm{E};\mathrm{F})\) 是一个 Lusin 空间（GT, IX, § 6, No.~7, 定理 3 的推论）。

推论. --- 在关于 E 的假设与定理 2 相同的前提下，再假设 E 的每个有界子集都是相对紧的。则 E 的强对偶是一个 Lusin 空间。* 特别地，满足第一可数公理的 Fréchet 空间若同时是 Montel 空间，则其强对偶是 Lusin 空间。*

* 例 2. --- 设 U 是 \(\mathbf{R}^{n}\) 的一个开子集。上述推论特别适用于 Fréchet 空间 \(E = \mathcal{C}^{\infty}(\mathrm{U})\) ；它的对偶 \(\mathcal{C}_{0}^{-\infty}(\mathrm{U})\) （U 上具有紧支集的广义函数所成的空间）于是是一个 Lusin 空间。空间 \(\mathcal{C}_{0}^{\infty}(\mathrm{U})\) 是满足第一可数公理的一列 Fréchet 空间 \(\mathcal{C}_{K_n}^{\infty}(\mathrm{U})\) 的严格归纳极限（III, 第 9 页）。我们可以证明，诸空间 \(\mathcal{C}_{K_n}^{\infty}(\mathrm{U})\) 中的每一个都是 Montel 空间；此外，\(\mathcal{C}_{0}^{\infty}(\mathrm{U})\) 的每个有界子集都含于诸空间 \(\mathcal{C}_{K_n}^{\infty}(\mathrm{U})\) 之一（III, 第 5 页, 命题 6）。于是可以应用定理 2 的推论。从而 \(\mathcal{C}_{0}^{-\infty}(\mathrm{U})\) （\(\mathcal{C}_{0}^{\infty}(\mathrm{U})\) 的对偶，U 上的广义函数空间）对强拓扑而言是一个 Lusin 空间。*

类似地可以证明：对 \(C^{n}\) 的每个开子集 U 以及 \(C^{n}\) 的每个紧子集 K，\(\mathcal{H}(U)\) 的强对偶与 \(\mathcal{H}(K)\) 的强对偶都是 Lusin 空间。*

注记. --- 设 E 如定理 2 中所示；设 F 是一个 Hausdorff 局部凸空间，它是一列连续线性映射 \(u_n: \mathrm{F}_n \to \mathrm{F}\) 的像的并，其中每个 \(\mathrm{F}_n\) 都是满足第一可数公理的 Fréchet 空间；则 \(\mathcal{L}_c(\mathrm{E};\mathrm{F})\) 是一个 Lusin 空间。如命题 1 中那样，我们先化归到每个 \(u_n\) 都是单射的情形；然后，如定理 2 的证明中那样，可以假设 E 是满足第一可数公理的 Fréchet 空间。于是由 I, 第 20 页, 命题 1，\(\mathcal{L}(\mathrm{E};\mathrm{F})\) 是诸 \(\mathcal{L}(\mathrm{E};\mathrm{F}_n)\) 的并；此外，典范单射 \(\mathcal{L}_c(\mathrm{E};\mathrm{F}_n) \to \mathcal{L}_c(\mathrm{E};\mathrm{F})\) 连续（GT, X, § 1, No.~4, 命题 3）。由于由定理 2，诸空间 \(\mathcal{L}_c(\mathrm{E};\mathrm{F}_n)\) 中的每一个都是 Lusin 空间，故 \(\mathcal{L}(\mathrm{E};\mathrm{F}_n)\) 对由 \(\mathcal{L}_c(\mathrm{E};\mathrm{F})\) 的拓扑诱导的拓扑而言也是 Lusin 空间（GT, IX, § 6, No.~4, 命题 11）；于是由 GT, IX, § 6, No.~7, 定理 3 的推论，\(\mathcal{L}_c(\mathrm{E};\mathrm{F})\) 是一个 Lusin 空间。

